\documentclass[12pt,a4paper]{article}
\usepackage{fontspec}
\usepackage{polyglossia}
\setmainlanguage{russian}
\setmainfont{Liberation Serif}
\usepackage[left=1.3cm,right=1.3cm,top=1.1cm,bottom=1.0cm]{geometry}
\usepackage{array}
\setlength{\parindent}{0pt}
\pagestyle{empty}
\renewcommand{\arraystretch}{1.45}

\begin{document}

\noindent{\Large\bfseries Обратный ход \ · \ ответы}\hfill{\small 15 сентября. Детям не выдаётся}
\vspace{4pt}\hrule\vspace{8pt}

\begin{tabular}{@{}c|p{7.6cm}|p{7.6cm}@{}}
\textbf{№} & \hfil\textbf{вариант A}\hfil & \hfil\textbf{вариант B}\hfil\\ \hline
1 & $10 \to 15 \to 5 \to 20 \to 14 \to 2$ \newline \textbf{задумала 10} & $9 \to 5 \to 15 \to 24 \to 4 \to 20$ \newline \textbf{задумал 9}\\ \hline
2 & $400 \to 800 \to 1600$ \newline \textbf{в 18:00} & $100 \to 200 \to 400 \to 800$ \newline \textbf{12 июля}\\ \hline
3 & $1600 \to 800 \to 400 \to 200 \to 100$, \newline пятый съел 50 \ \textbf{было 1600 граммов} & $960 \to 480 \to 240 \to 120 \to 60$, \newline пятый съел 30 \ \textbf{было 960 граммов}\\ \hline
4 & \textbf{14-й день} & \textbf{18-й день}\\ \hline
5 & $21 \to 18 \to 12 \to 0$ \newline \textbf{21 монета} & $28 \to 24 \to 16 \to 0$ \newline \textbf{28 монет}\\ \hline
6 & $(39,21,12) \to (6,42,24) \to$ \newline $(12,12,48) \to (24,24,24)$ \newline \textbf{Аня 39, Боря 21, Валя 12} & $(52,28,16) \to (8,56,32) \to$ \newline $(16,16,64) \to (32,32,32)$ \newline \textbf{Аня 52, Боря 28, Валя 16}\\ \hline
7 & $127 \to 63 \to 31 \to 15 \to 7 \to 3 \to 1 \to 0$ \newline \textbf{127 гусей} & $63 \to 31 \to 15 \to 7 \to 3 \to 1 \to 0$ \newline \textbf{63 гуся}\\ \hline
8$^*$ & $458 \to 45 \to 90 \to 9 \to 18 \to 36 \to$ \newline $72 \to 144 \to 14$ \newline \textbf{годится любая верная цепочка} & $926 \to 92 \to 9 \to 18 \to 36 \to 72 \to$ \newline $144 \to 14$ \newline \textbf{годится любая верная цепочка}\\ \hline
9$^*$ & $23 \to 7 \to 1 \to 0$ \newline \textbf{23 задачи} & $35 \to 11 \to 2 \to 0$ \newline \textbf{35 задач}\\ \hline
10$^*$ & назад из 74: $74,\,47,\,37,\,73$ \newline назад из 112: $112,\,121,\,211,\,56,\,65,$ \newline $28,\,82,\,41,\,14,\,7$ \newline \textbf{нельзя ни то, ни другое} & назад из 38: $38,\,83,\,19,\,91$ \newline назад из 228: $228,\,282,\,822,\,114,$ \newline $141,\,411,\,57,\,75$ \newline \textbf{нельзя ни то, ни другое}\\
\end{tabular}

\vspace{14pt}\hrule\vspace{8pt}

\textbf{Задача 6.} Сумма фишек не меняется: 72 в варианте A, 96 в B.

\vspace{3pt}
\textbf{Задача 8.} Цепочка засчитывается любая, лишь бы каждый переход был удвоением
или стиранием последней цифры. Пересчитать глазами.

\vspace{3pt}
\textbf{Задача 10.} Обе цепочки назад конечны, и единицы в них нет. Перестановка цифр
не меняет их количества, а поделить пополам можно только чётное --- поэтому список
замыкается сам на себе.

\vspace{12pt}\hrule\vspace{5pt}
{\small Все цепочки пересчитаны программой 2026-09-14, а не переписаны из источников.}

\end{document}
